\documentclass{article}
\usepackage[a4paper, total={6.5in, 10in}]{geometry}
\setlength{\parindent}{0pt}
\usepackage{tcolorbox}
\tcbset{colback=white!94!black, colframe=black!60!white, coltitle = white, boxrule=0.8pt, arc=2mm}
\newtcolorbox{boxs}[2][]{title= #2,#1}
\usepackage{datetime2}
\usepackage[british]{babel}
\usepackage{amsfonts}
\usepackage{amsmath}

\title{\textbf{Limits}}
\author{Robert O'Connell}
\date{\today}
\begin{document}
\maketitle
\subsection*{Epsilon-Delta Definition of Limits}
We say that \(f(x)\) approaches the limit \(L\) as \(x\) approaches \(x_0\) and we write
\[
\lim_{x\to x_0}{f(x)} = L
\]
if given any \(\epsilon > 0 \) there exists \(\delta > 0 \) such that
\[
0 \neq |x-x_0| < \delta \rightarrow |f(x)-L| < \epsilon
\]
(for examples see written notes)

\subsection*{Triangle Inequality}
One has \(|x+y| \leq |x| + |y|\) for all \(x,y\)

For real numbers this holds because
\begin{align*}
    |x+y|^2 &= (x+y)^2 \\
    &= x^2 +y^2+2xy \\
    &\leq |x| + |y| + 2|x||y| = (|x| + |y|)^2
\end{align*}


\subsection*{Theorem - Properties of limits}
The limit of a sum/product/quotient is the sum/product/quotient of the limits

Namely, if
\[
\lim_{x\to x_0}{f(x)} = L, \quad \lim_{x\to x_0}{g(x)} = M
\]
Then
\[
\lim_{x\to x_0}{[f(x)+g(x)]} = L + M
\]
\[
\lim_{x\to x_0}{[f(x)\cdot g(x)]} = L \cdot M
\]
\[
\lim_{x\to x_0}{\left[\frac{f(x)}{g(x)}\right]} = \frac{L}{M}, \quad M \neq 0
\]

\subsubsection*{Proofs}
We use the \(\epsilon - \delta\) definition

We know \(\lim_{x\to x_0}{f(x)} = L\) and \(\lim_{x\to x_0}{g(x)} = M\)
\\

\textbf{To compute the limit of the sums}

We need \(\lim_{x\to x_0}{[f(x) + g(x)]} = L + M\)

Now, let \(\epsilon > 0 \) be given, we estimate
\begin{align*}
    |f(x) + g(x) - (L+M)| &= |f(x) - L + g(x)-M| \\
    &\leq |f(x)-L| + |g(x)-M| 
\end{align*}

Since \(\lim_{x\to x_0}{f(x)} = L\) there exists \(\delta_1 > 0 \)
\[
0 \neq |x-x_0|<\delta_1 \Rightarrow |f(x)-L| < \frac{\epsilon}{2}
\]

Since \(\lim_{x\to x_0}{g(x)} = M\) there exists \(\delta_2 > 0 \)
\[
0 \neq |x-x_0|<\delta_2 \Rightarrow |g(x)-M| < \frac{\epsilon}{2}
\]

Thus,
\[
0 \neq |x-x_0| < \text{min}\{\delta_1,\delta_2\} \Rightarrow |f(x) + g(x)-L-M| < \epsilon
\]
\\

\textbf{To compute the limit of the prodcuts}

We estimate,
\begin{align*}
    |f(x)\cdot g(x)-LM| &= |f(x)\cdot g(x) -L\cdot g(x) + L \cdot g(x) -LM| \\
    &\leq |g(x)||f(x) - L| + |L||g(x)-M|
\end{align*}
\end{document}